\color{uniblue}\section[Техническое обслуживание и ремонт]{Техническое обслуживание и ремонт}\label{sec:techrepair}
\color{black}

\color{uniblue}\subsection{Техническое обслуживание устройства}\color{black}

\begin{enumerate}[label=\arabic{section}.\arabic{subsection}.\arabic*, labelsep=4pt, leftmargin=0pt, itemindent=57pt]

\item\sloppy
Процесс  технического  обслуживания  устройств  релейной  защиты  регламентирован  приказом Минэнерго  России  «Об  утверждении  Правил  технического  обслуживания  устройств  и  комплексов релейной  защиты  и  автоматики  и  внесении  изменений  в  требования  к  обеспечению  надежности электроэнергетических  систем,  надежности  и  безопасности  объектов  электроэнергетики  и энергопринимающих установок».
\item
Техническое обслуживание устройств <<ЮНИТ-М300>> должно выполняться в соответствии с документом ЮТКБ.656122.609 ИС – <<Устройства защиты и автоматики серии <<ЮНИТ-М300>>. Инструкция эксплуатационная специальная. Методические указания по выполнению технического обслуживания>>.

\end{enumerate}

\color{uniblue}\subsection{Текущий ремонт}\color{black}

\begin{enumerate}[label=\arabic{section}.\arabic{subsection}.\arabic*, labelsep=4pt, leftmargin=0pt, itemindent=57pt]

\item
Устройство серии <<ЮНИТ-М300>>, как было описано выше, выполнено с использованием микропроцессорных технологий, поэтому необходимость в его текущем ремонте отсутствует. Работоспособность и безотказность работы устройства обеспечивается функционированием системы самодиагностики.
\item
Углубленная самодиагностика аппаратных средств, а также состояние программной памяти, памяти данных  и  буферной  памяти  производится  при  включении  устройства.  Поэтому  при  возникновении сообщений о критических или некритических ошибках от системы самодиагностики, даже при возможности устранения  их  своими  силами,  следует  обратиться  на  предприятие-изготовитель для  консультации или проведения регламентных работ.
\item 
Поводом  для  обращения  на  предприятие-изготовитель  являются  также  признаки  нехарактерного нагрева или перегрева отдельных элементов, участков или частей устройства.
\item
Вышедшие  из  строя  устройства  должны  проходить  ремонт  на  предприятии-изготовителе, обеспечивающем  гарантийное  и  послегарантийное  обслуживание, адрес которого  указан  в  паспорте устройства.  В  качестве  ЗИП,  если  это  предусмотрено  договором  поставки,  могут  предоставляться устройства  <<ЮНИТ-М300>> с необходимыми  конфигурациями,  позволяющие  оперативно  произвести  замену вышедших из строя устройств.
\item
Перечень сигналов самодиагностики приведен в документации ЮТКБ.656122.609 РЭ <<Устройство защиты и автоматики серии <<ЮНИТ-М300>>. Руководство по эксплуатации общее>>.
\end{enumerate}

